\documentclass[a4paper,11pt]{article}
\usepackage[utf8]{inputenc}
\usepackage[T1]{fontenc}
\usepackage[french]{babel}
\usepackage{lmodern}
\usepackage{amsmath,amssymb,amsthm}
\usepackage{mathrsfs}
\usepackage{enumitem}
\usepackage{multicol}

\setlist[enumerate]{label=\textbf{\textcolor{Couleur8}{\arabic*.}}, left=1em}

% Si vous voulez aussi modifier les listes enumerate imbriquées :
\setlist[enumerate,2]{label=\textcolor{Couleur8}{\alph*)}}
\setlist[enumerate,3]{label=\textcolor{Couleur8}{\roman*)}}
\setlist[enumerate,4]{label=\textcolor{Couleur8}{\Alph*)}}
\usepackage{graphicx}
\usepackage{xcolor}
\usepackage{tcolorbox}
\usepackage{geometry}
\usepackage{fancyhdr}
\usepackage{titlesec}
\usepackage{cancel}
\usepackage{ifthen}
\usepackage{tikz}
\usetikzlibrary{arrows.meta}
\usepackage{tkz-tab}
\usepackage{eurosym}

% OPTION PROF/ÉLÈVE
% Décommentez UNE des deux lignes suivantes :
\newboolean{prof}\setboolean{prof}{true}   % Version professeur (avec corrections des devoirs)
\newboolean{prof}\setboolean{prof}{true}  % Version élève (sans corrections des devoirs)

\definecolor{Couleur1}{HTML}{4A5D7A} %Th
\definecolor{Couleur2}{HTML}{A5B5D6} %Prop
\definecolor{Couleur3}{HTML}{A8C68A} %Méthode
\definecolor{Couleur6}{HTML}{8A9CC1} %Def
\definecolor{Couleur7}{HTML}{E6B459} %Rq Voc
\definecolor{Couleur8}{HTML}{D36740} %Titres
\definecolor{correction}{HTML}{D36740} % Couleur pour les corrections
\definecolor{coopmathscorr}{HTML}{F15929} % Couleur corrections coopmaths

% Configuration des marges
\geometry{
	top=2cm,
	bottom=2cm,
	inner=1.5cm,
	outer=1.5cm,
}

% Police
\renewcommand{\familydefault}{\sfdefault}

% Compteurs
\newcounter{seancenum}
\newcounter{seancenumD}
\setcounter{seancenum}{1}
\setcounter{seancenumD}{1}

% Configuration des en-têtes et pieds de page
\pagestyle{fancy}
\fancyhf{}
\ifthenelse{\boolean{prof}}{
    \fancyhead[L]{\textcolor{Couleur8}{\textbf{Corrections Automatismes - Classe de seconde - Version Professeur}}}
}{
    \fancyhead[L]{\textcolor{Couleur8}{\textbf{Corrections Devoirs - Séance 28 - Classe de seconde}}}
}
\fancyhead[R]{\textcolor{Couleur8}{\textbf{2025-2026}}}
\fancyfoot[L]{\textcolor{Couleur8}{Mme Bertail et Mme Pinto}}
\fancyfoot[R]{\textcolor{Couleur8}{\thepage}}
\renewcommand{\headrulewidth}{1pt}
\renewcommand{\headrule}{\hbox to\headwidth{\color{Couleur8}\leaders\hrule height \headrulewidth\hfill}}

% Environnements personnalisés
\newtcolorbox{seance}[1]{
    colback=Couleur6!10,
    colframe=Couleur1!65,
    fonttitle=\bfseries\large,
    title=Correction Séance \theseancenum #1,
    left=5mm,
    right=5mm,
    top=3mm,
    bottom=3mm,
    arc=0mm,
    after=\stepcounter{seancenum}
}

\newtcolorbox{seanceD}[1]{
    colback=Couleur6!5,
    colframe=Couleur6!45,
    fonttitle=\bfseries\large,
    title=Correction Devoirs - Séance \theseancenumD #1,
    left=5mm,
    right=5mm,
    top=3mm,
    bottom=3mm,
    arc=0mm,
    after=\stepcounter{seancenumD}
}

\begin{document}

\setcounter{seancenumD}{28}
\newpage
\begin{seanceD}{} %28

\begin{enumerate}
\item La série triée par ordre croissant est : $2  ;  6  ;  7  ;  10  ;  16  ;  30$.\\La série contient $6$ valeurs.\\
      Pour trouver le rang de $Q_3$, on calcule les trois quarts de $6$ qui vaut
       $\dfrac{3\times6}{4}=4{,}5$.\\On arrondit à l'entier supérieur qui vaut $5$ .\\ Le troisième quartile est donc la valeur de rang $5$ de la série classée : $Q_3=16$.\\La bonne réponse est la réponse {\bfseries \color[HTML]{F15929}E}.

\item La somme des probabilités doit être égale à 1.  \\Comme $0{,}3=\dfrac{3}{10}$, on a : \\
        $\begin{aligned}
        x&=1-\left(\dfrac{1}{6}+\dfrac{3}{10}+\dfrac{1}{5}\right)\\
        x&=1-\left(\dfrac{5}{30}+\dfrac{9}{30}+\dfrac{6}{30}\right)\\
        x&=1-\dfrac{2}{3}\\
        x&={\color[HTML]{F15929}\boldsymbol{\dfrac{1}{3}}}
        \end{aligned}$ 

\medskip
La bonne réponse est la réponse {\bfseries \color[HTML]{F15929}D}.

\item On reconnaît la droite grâce à son ordonnée à l'origine ($-1<0$) et son coefficient directeur ($1>0$).\\
    Il s'agit de la droite coupant l'axe des ordonnées en-dessous de l'axe des abscisses et qui monte.\\
    Il s'agit de la droite ${\color[HTML]{F15929}\boldsymbol{D_2}}$.\\La bonne réponse est la réponse {\bfseries \color[HTML]{F15929}B}.

\item À partir des évolutions en pourcentage, on déduit les coefficients multiplicateurs : \\On note $CM_1 = 1 + \dfrac{10}{100}=1{,}1$ et $CM_2 = 1 + \dfrac{40}{100}=1{,}4$.\\ Le coefficient multiplicateur global est : \\ $CM = CM_1 \times CM_2 = 1{,}1 \times 1{,}4 = 1{,}54$ \\Or, multiplier par $1{,}54$ revient à avoir {\bfseries \color[HTML]{F15929}une} {\bfseries \color[HTML]{F15929}augmentation} {\bfseries \color[HTML]{F15929}de} ${\color[HTML]{F15929}\boldsymbol{54\,\%}}$.\\La bonne réponse est la réponse {\bfseries \color[HTML]{F15929}A}.

\item L'équation $-9x^2-8x+1=1$ s'écrit $-9x^2-8x=0$.\\
          En factorisant le premier membre (facteur commun $x$), on obtient $x(-9x-8)=0$.\\
          On reconnaît une équation produit nul dont les solutions sont : $0$ et $\dfrac{8}{-9}=-\dfrac{8}{9}$.\\
          $\mathscr{S}={\color[HTML]{F15929}\boldsymbol{\{0;-\dfrac{8}{9}\}}}$\\La bonne réponse est la réponse {\bfseries \color[HTML]{F15929}C}.

\item De la relation $-7x-6c=2$, on déduit en ajoutant $7x$ dans chaque membre :
          $-6c=2+7x$\\ Puis, en divisant par $-6$, on obtient : $c=\dfrac{2+7x}{-6}$ que l'on peut écrire également : ${\color[HTML]{F15929}\boldsymbol{ c=\dfrac{-2-7x}{6}}}$.\\La bonne réponse est la réponse {\bfseries \color[HTML]{F15929}A}.

\end{enumerate}

\end{seanceD}

\end{document}